% !TeX root = ../../main.tex
\subsection{Low-complexity vertex integrals}
\label{subsec:a:03}

The four small vertex functions needed below follow from the pure-psi
string and dilaton equations and the forgetting formulas of
Subsection~\ref{subsec:s02:04}. We use the stack initial values
\[
 \int_{\MM_{0,3}}1=1,\qquad
 \int_{\MM_{1,1}}\psi_1=\frac1{24}.
\]
In the next proof all displayed integrals without an exponential are
ordinary, unnormalized intersection numbers. In particular they are
neither the Introduction's Hodge bracket nor the normalized WP brackets.

\begin{lmm}[Small vertex polynomials]\label{lmm:a:small-vertices}
For complex cotangent parameters of the indicated lengths,
\begin{align}
 W^{\rm WP}_{0,3}(t_1,t_2,t_3)&=1,\notag\\
 W^{\rm WP}_{1,1}(t)&=\frac{\pi^2}{12}-\frac{t}{24},\notag\\
 W^{\rm WP}_{0,4}(t_1,t_2,t_3,t_4)&=2\pi^2-\sum_{i=1}^4t_i,\notag\\
 W^{\rm WP}_{1,2}(t_1,t_2)&=\frac{\pi^4}{4}
       -\frac{\pi^2}{6}(t_1+t_2)+\frac{(t_1+t_2)^2}{48}.
 \label{eq:a:small-vertices}
\end{align}
In particular the complete stack length polynomial is
$V_{1,1}(y)=\pi^2/12+y^2/48$.
\end{lmm}
\begin{prf}
For nonnegative integer exponents, string removes a zero-degree marking
by summing over decrements of the other positive degrees, and dilaton
removes a degree-one marking with factor $2G-2+n$, where $n$ is the
number of remaining markings. Together with the two initial values,
these give
\[
 \int_{\MM_{0,4}}\psi_i=1,\qquad
 \int_{\MM_{1,2}}\psi_1^2
 =\int_{\MM_{1,2}}\psi_1\psi_2=\frac1{24}.
\]
The one-point forgetting definition of $\kappa_1$ gives
\[
 \int_{\MM_{0,4}}\kappa_1
       =\int_{\MM_{0,5}}\psi_5^2=1,\qquad
 \int_{\MM_{1,1}}\kappa_1
       =\int_{\MM_{1,2}}\psi_2^2=\frac1{24}.
\]
For the mixed integral and the square of $\kappa_1$ the same equations
and the two-input inversion~\eqref{eq:s02:partition-inversion} give
\begin{align*}
 \int_{\MM_{1,2}}\kappa_1\psi_1
  &=\int_{\MM_{1,3}}\psi_1\psi_3^2=\frac1{12},\\
 \int_{\MM_{1,2}}\kappa_1^2
  &=\int_{\MM_{1,4}}\psi_3^2\psi_4^2
                -\int_{\MM_{1,3}}\psi_3^3
    =\frac16-\frac1{24}=\frac18.
\end{align*}
All new-point powers in these formulas are at least two. Therefore the
rational-tail lemma of Section~\ref{subsec:s02:04} permits replacing
old pulled-back cotangent classes by source classes. The minus sign
in the last line is the merged-block term of the inversion, and must
be retained.

Expanding $\exp(\alpha\kappa_1-\sum t_i\psi_i)$ now proves each
polynomial, including its factorials. For instance its constant term
on $\MM_{1,2}$ is $(\alpha^2/2)(1/8)=\pi^4/4$ and its linear
coefficient is $-\alpha/12=-\pi^2/6$.
Substituting $t=-y^2/2$ gives the stated length polynomial.
\end{prf}

\begin{rmk}[Conventions]
The initial values, string and dilaton appear in
\cite[Section~4.1 and (B.1)]{W26}. The length interpretation is
\cite[Theorem~4.4]{Mir07}, with the stack convention explained in
Section~\ref{subsec:s03:01}. The entire $(1,1)$ polynomial, including
its constant term, is half the effective geometric polynomial in
Mirzakhani's convention. We never halve only its descendant term.
\end{rmk}

The normalized single-degree brackets used for the retained small
components are as follows.
Entries beyond the dimension are zero:
\begin{equation}\label{eq:a:small-amplitudes}
\begin{array}{c|ccc}
 (G,n)&[\tau_0^n]_{G,n}&[\tau_1\tau_0^{n-1}]_{G,n}
                                      &[\tau_2\tau_0^{n-1}]_{G,n}\\\hline
 (0,3)&1&0&0\\
 (1,1)&\pi^2/12&1/2&0\\
 (0,4)&2\pi^2&12&0\\
 (1,2)&\pi^4/4&2\pi^2&10
\end{array}
\end{equation}
They follow directly from~\eqref{eq:s03:descendants}; no volume
asymptotic is used to obtain a low-type amplitude.
